%! Trigonometry and Polar Coordinates — Unit 3, Lesson 3.13 — Group Activity
\documentclass[10pt]{article}
\usepackage{precalculus-article}
\usepackage{precalculus-boxes}
\newcommand{\TallMath}[1]{$\displaystyle #1\rule[-1.4em]{0pt}{3.2em}$}

\begin{document}

\pageheader{Unit 3, Lesson 3.13}{Group Activity}
\namepartnerperiod

\vspace{0.10in}

%----------------------------------------------------------------------
% Context
%----------------------------------------------------------------------
\begin{scenariobox}[Context: Polar Coordinates in Navigation and Radar]{sky}
Air traffic control radar systems report the location of aircraft using a
\emph{distance} (range) and a \emph{bearing} (angle from north) --- not $x$
and $y$ coordinates. This is the polar coordinate system in action. In this
activity your group will practice converting between polar and rectangular
representations and explore the fact that a single location has many valid
polar descriptions. Choose \textbf{one tier} based on your readiness level.
Everyone should be prepared to explain their reasoning.
\end{scenariobox}

\vspace{0.08in}

%----------------------------------------------------------------------
% Tier R
%----------------------------------------------------------------------
\begin{tcolorbox}[colback=white, colframe=black!40,
    title=\textbf{Tier R --- Reinforce},
    fonttitle=\bfseries, arc=2mm, left=3mm, right=3mm, top=2mm, bottom=2mm]

\textbf{Directions:} For each polar coordinate pair, find the rectangular
coordinates using $x = r\cos\theta$ and $y = r\sin\theta$. Show each step.

\vspace{6pt}
\renewcommand{\arraystretch}{2.2}
\begin{tabular}{|c|c|c|c|c|}
\hline
\textbf{Polar point} & \textbf{$r$} & \textbf{$\theta$} & \textbf{$x = r\cos\theta$} & \textbf{$y = r\sin\theta$} \\
\hline
$(4,\, \pi/6)$ & 4 & $\pi/6$ & \underline{\hspace{2.8cm}} & \underline{\hspace{2.8cm}} \\
\hline
$(2,\, \pi/2)$ & 2 & $\pi/2$ & \underline{\hspace{2.8cm}} & \underline{\hspace{2.8cm}} \\
\hline
$(6,\, \pi)$ & 6 & $\pi$ & \underline{\hspace{2.8cm}} & \underline{\hspace{2.8cm}} \\
\hline
$(3,\, 3\pi/2)$ & 3 & $3\pi/2$ & \underline{\hspace{2.8cm}} & \underline{\hspace{2.8cm}} \\
\hline
\end{tabular}

\vspace{8pt}
\textbf{1.}\quad The point $(4, \pi/6)$ gives rectangular coordinates
$(x,y) = (2\sqrt{3}, 2)$. Verify by computing $\sqrt{x^2+y^2}$:

\vspace{4pt}
$\sqrt{(2\sqrt{3})^2 + 2^2} = \sqrt{\underline{\hspace{1cm}} + \underline{\hspace{1cm}}} = \sqrt{\underline{\hspace{1cm}}} = $ \underline{\hspace{1.5cm}}

Does this equal the original $r$? \underline{\hspace{1.2cm}}

\vspace{8pt}
\textbf{2.}\quad The polar point $(2, \pi/2)$ lies on the positive $y$-axis.
Explain why $x = 0$ in this case without computing $2\cos(\pi/2)$.

\writelines{2}

\vspace{6pt}
\textbf{3.}\quad Without computing, predict the rectangular coordinates
of the polar point $(5, 0)$. Explain your reasoning.

\vspace{4pt}
Rectangular: $(\underline{\hspace{1.5cm}},\, \underline{\hspace{1.5cm}})$ \quad Reason: \writeline

\end{tcolorbox}

\vspace{0.08in}

%----------------------------------------------------------------------
% Tier A
%----------------------------------------------------------------------
\begin{tcolorbox}[colback=white, colframe=black!40,
    title=\textbf{Tier A --- Apply},
    fonttitle=\bfseries, arc=2mm, left=3mm, right=3mm, top=2mm, bottom=2mm]

\textbf{Directions:} Complete all three parts. Show all work and quadrant reasoning.

\vspace{6pt}
\textbf{Part 1 --- Rectangular $\to$ Polar} ($r > 0$, $0 \le \theta < 2\pi$)

\vspace{4pt}
(a)\quad $(1,\, \sqrt{3})$: \quad $r = $ \underline{\hspace{1.8cm}} \quad
$\theta = $ \underline{\hspace{1.8cm}} \quad Polar: \underline{\hspace{2.5cm}}

\vspace{6pt}
(b)\quad $(-4,\, 0)$: \quad $r = $ \underline{\hspace{1.8cm}} \quad
$\theta = $ \underline{\hspace{1.8cm}} \quad Polar: \underline{\hspace{2.5cm}}

\vspace{6pt}
(c)\quad $(3,\, -3)$: \quad $r = $ \underline{\hspace{1.8cm}} \quad
$\theta = $ \underline{\hspace{1.8cm}} \quad Polar: \underline{\hspace{2.5cm}}

\vspace{10pt}
\textbf{Part 2 --- Multiple Representations}

Given the polar point $(2,\, \pi/3)$, write two other polar representations
of the same point.

\vspace{4pt}
Representation 1 (add $2\pi$): \underline{\hspace{5cm}}

\vspace{4pt}
Representation 2 (use $-r$): \underline{\hspace{5cm}}

\vspace{4pt}
Verify Representation 2 by converting it to rectangular and comparing with
$(2,\, \pi/3)$ in rectangular form.

\vspace{4pt}
$(2,\pi/3)$ in rectangular: $(\underline{\hspace{1.5cm}},\, \underline{\hspace{1.5cm}})$
\quad Rep.\ 2 in rectangular: $(\underline{\hspace{1.5cm}},\, \underline{\hspace{1.5cm}})$
\quad Match? \underline{\hspace{1cm}}

\vspace{10pt}
\textbf{Part 3 --- Complex Number in Polar Form}

Write $2 - 2i$ in polar form. Show all steps.

\vspace{4pt}
$r = \sqrt{(\underline{\hspace{0.8cm}})^2 + (\underline{\hspace{0.8cm}})^2} = $ \underline{\hspace{2cm}}

\vspace{4pt}
Quadrant of $2-2i$: \underline{\hspace{1.5cm}} \qquad
$\theta = \arctan\!\left(\dfrac{-2}{2}\right) = $ \underline{\hspace{2cm}}
(adjust if needed: $\theta = $ \underline{\hspace{2cm}})

\vspace{4pt}
Polar form: \underline{\hspace{6cm}}

\end{tcolorbox}

\vspace{0.08in}

%----------------------------------------------------------------------
% Tier E
%----------------------------------------------------------------------
\begin{tcolorbox}[colback=white, colframe=black!40,
    title=\textbf{Tier E --- Extend},
    fonttitle=\bfseries, arc=2mm, left=3mm, right=3mm, top=2mm, bottom=2mm]

\textbf{Directions:} Complete both problems. Justify your reasoning carefully.

\vspace{6pt}
\textbf{Problem 1 --- General Polar Representations of a Q2 Point}

Let $(x,y)$ be a point in Quadrant 2 with $r = \sqrt{x^2+y^2}$ and
reference angle $\alpha = \arctan(|y/x|)$.

\vspace{4pt}
(a)\quad Write the standard polar form of the point in Q2 using $r$ and $\alpha$.

\vspace{4pt}
$(r,\theta) = (\underline{\hspace{1.5cm}},\; \underline{\hspace{3.5cm}})$

\vspace{6pt}
(b)\quad Write the general formula for \emph{all} polar representations
$(r,\theta)$ of this Q2 point (let $k$ be any integer).

\vspace{4pt}
All representations: $(\underline{\hspace{1.5cm}},\; \underline{\hspace{4cm}})$ for any integer $k$,
or $(\underline{\hspace{1.5cm}},\; \underline{\hspace{4cm}})$ for any integer $k$.

\vspace{8pt}
\textbf{Problem 2 --- Showing $(-r,\theta)$ and $(r,\theta+\pi)$ Name the Same Point}

Two polar pairs $(-r,\theta)$ and $(r,\theta+\pi)$ are claimed to represent the
same rectangular point. Prove this algebraically.

\vspace{4pt}
Rectangular coordinates of $(-r,\theta)$:
\begin{align*}
x_1 &= (-r)\cos\theta = \underline{\hspace{4cm}} \\
y_1 &= (-r)\sin\theta = \underline{\hspace{4cm}}
\end{align*}

Rectangular coordinates of $(r,\theta+\pi)$:
\begin{align*}
x_2 &= r\cos(\theta+\pi) = r \cdot (\underline{\hspace{3cm}}) = \underline{\hspace{3cm}} \\
y_2 &= r\sin(\theta+\pi) = r \cdot (\underline{\hspace{3cm}}) = \underline{\hspace{3cm}}
\end{align*}

Therefore $x_1 = x_2$ and $y_1 = y_2$, so the two representations name the
\underline{\hspace{2cm}} rectangular point. \qquad $\square$

\vspace{6pt}
(c)\quad Why is this result important for the uniqueness of polar coordinates?

\writelines{2}

\end{tcolorbox}

\end{document}
